\chapter*{נספח א': מילון מונחים}
\addcontentsline{toc}{chapter}{נספח א': מילון מונחים}

נספח זה מרכז את המונחים המרכזיים המופיעים בעבודה, תוך מתן הגדרות תמציתיות לשימוש עתידי.

\begin{description}

\item[\textenglish{Attention Mechanism} (מנגנון תשומת לב)]
תהליך חישובי המאפשר לכל אסימון (\textenglish{token}) בנתון קלט לקבוע את מידת הרלוונטיות של יתר האסימונים בייצוג התוצאה. מחושב באמצעות מטריצות שאילתה (\textenglish{Query}), מפתח (\textenglish{Key}) וערך (\textenglish{Value}).

\item[\textenglish{Transformer}]
ארכיטקטורת רשת נוירונית המבוססת אך ורק על מנגנון תשומת לב, ללא רשתות רקורנטיות. הוצגה לראשונה במאמר \textenglish{"Attention Is All You Need"} (2017) והפכה לבסיס של רוב מודלי השפה המודרניים.

\item[\textenglish{Large Language Model} (\textenglish{LLM})]
מודל שפה גדול המאומן על כמויות עצומות של טקסט, המסוגל לבצע מגוון רחב של משימות שפה ללא כוונון ייעודי. דוגמאות: \textenglish{GPT-4}, \textenglish{Claude}, \textenglish{Gemini}.

\item[\textenglish{Model Context Protocol} (\textenglish{MCP})]
תקן פתוח המגדיר ממשק אחיד בין מודלי שפה לכלים חיצוניים. מאפשר לסוכן לגלות, לפנות ולהשתמש בכלים (שרתים) בצורה עקבית ובטוחה.

\item[\textenglish{ReAct} (חשיבה-פעולה)]
מסגרת לסוכני \textenglish{LLM} המשלבת שלבי חשיבה מפורשים (\textenglish{Reasoning}) עם שלבי פעולה (\textenglish{Acting}) בלולאה חוזרת. מאפשרת לסוכן לנמק לפני שהוא מבצע קריאה לכלי.

\item[\textenglish{Reflexion}]
שיטה לשיפור עצמי של סוכנים על בסיס משוב מילולי. הסוכן מסכם את כישלונותיו הקודמים ומאחסן אותם בזיכרון, ומשתמש בהם לשיפור ניסיונות עתידיים.

\item[\textenglish{Chain-of-Thought} (\textenglish{CoT})]
טכניקת הנחיה המעודדת את המודל לפרט את שלבי החשיבה הביניים לפני מתן התשובה הסופית, ובכך משפרת ביצועים במשימות מורכבות.

\item[\textenglish{Toolformer}]
מודל שפה שאומן לשלב באופן אוטונומי קריאות \textenglish{API} בתוך יצור הטקסט שלו. המודל לומד מתי כדאי לקרוא לכלי ואיך לפרש את תוצאותיו.

\item[\textenglish{AvaTaR}]
מסגרת רב-סוכנית המשתמשת בלמידה מחיזוק להכשרת סוכני כלים. בכל איטרציה, סוכן מתאמן (\textenglish{optimizer}) משפר את מדיניות השימוש בכלים של הסוכן הראשי.

\item[\textenglish{DEPS} (תכנון מודע-תלויות)]
שיטת תכנון המגדירה גרף תלויות בין תת-משימות ומבטיחה ביצוע בסדר הנכון, תוך ניהול תלויות נתונים בין שלבים.

\item[\textenglish{BiDi} (דו-כיווניות)]
טכנולוגיה לעיבוד טקסט המשלב כיוונים שונים: ימין-לשמאל (עברית) ושמאל-לימין (אנגלית, נוסחאות). מוטמעת בהפקת המסמך באמצעות חבילות \textenglish{polyglossia} ו-\textenglish{luabidi}.

\item[\textenglish{Fine-Tuning} (כוונון עדין)]
תהליך אימון נוסף של מודל קיים על מערך נתונים ממוקד, במטרה להתאים את התנהגותו למשימה ספציפית מבלי לאמן מאפס.

\item[\textenglish{Hallucination} (הזיה)]
תופעה שבה מודל שפה מייצר מידע שגוי שנראה סביר, אך אינו מבוסס על נתוני האימון או על המציאות. אחד האתגרים המרכזיים בפריסת \textenglish{LLM} בסביבות ייצוריות.

\end{description}
